\documentclass[reprint,twocolumn,prb,amsmath,amssymb,superscriptaddress,floatfix]{revtex4-2}

\usepackage{graphicx}
\usepackage{bm}
\usepackage[T1]{fontenc}
% Times-compatible text and math, matching the STIX serif used in the figures,
% so that a symbol in a caption and the same symbol inside a panel agree.
\usepackage{newtxtext}
\usepackage{newtxmath}
\usepackage{booktabs}
\usepackage{xcolor}
\usepackage{xspace}

% ---------------------------------------------------------------------------
% Pastel link inks.  Soft enough to sit quietly in 10pt text, saturated enough
% that a link never reads as plain type.  Three hues separate the kind of link:
%   pastellink  desaturated blue    internal cross-references (Eq., Fig., Sec.)
%   pastelcite  desaturated teal    bibliography citations
%   pasteurl    desaturated purple  URLs, DOIs, file links
% ---------------------------------------------------------------------------
\definecolor{pastellink}{rgb}{0.35,0.45,0.75}
\definecolor{pastelcite}{rgb}{0.20,0.52,0.45}
\definecolor{pastelurl}{rgb}{0.52,0.36,0.60}

% hyperref last (RevTeX 4-2 supports it natively).
\usepackage[colorlinks=true,
            linkcolor=pastellink,
            citecolor=pastelcite,
            urlcolor=pastelurl,
            filecolor=pastelurl,
            menucolor=pastellink,
            anchorcolor=pastellink,
            linktoc=all,
            breaklinks=true]{hyperref}

\graphicspath{{figures/}}

% Every number still owed by the running DFT+U campaign lives here.
% ---------------------------------------------------------------------------
% dft_numbers.tex
%
% Every number in the manuscript that comes from the DFT+U campaign is defined
% here and nowhere else, in both the article and the Supplemental Material.
%
%   campaign (a)  results/dftu_row_ordered      row-ordered Li_{1/2}CoO2
%   campaign (b)  results/dftu_order_contrast   row / zigzag / clumped, x = 1/2
%   campaign (c)  results/dftu_x13              sqrt3 x sqrt3 Li_{1/3}CoO2
%   distillation  results/dftu_summary.json
%
% Usage convention: a macro is followed by "~" and then its unit, e.g.
% \comomentRow~$\mu_{\mathrm{B}}$.  "~" is in the xspace exception list.
%
% PROVENANCE NOTE (verification finding B7).  The x = 1/2 site moments, charge
% split and the moment ceiling derived from them are quoted from the ground
% state, afmlo_intra_U4.0 in cell B1, and not from the ferromagnetic solution of
% the seven-atom cell.  \comomentRowRange records the spread across the FM and
% both AFM configurations, which is the honest measure of how well determined
% the moment is.
%
% RETIRED MACROS, and why.  \deltaEOrder and the original reading of
% \dosContrast assumed the row motif lies below the zigzag motif and the ordered
% motif carries the smaller density of states at E_F; neither survived.
% \dosContrast now names the measured U = 0 ratio, which is the falsification
% itself.  \momentCellRow, \enmfmRow, \dosRow, \dosZigzag, \dosClumped,
% \dosXthird and \deltaMdft are retired, the last because the bulk ground states
% are antiferromagnetic and supply no net Delta M.  \jintraRow and \jinterRow
% are retired in favour of \eafmIntra and \eafmInter, which are signed energy
% differences and carry no implicit spin-Hamiltonian convention; \jeffXthird is
% retired in favour of \splitXthird, the raw splitting per cell.
% ---------------------------------------------------------------------------

\newcommand{\Uused}{\ensuremath{4.0}\xspace}

% --- (a) row-ordered Li_{1/2}CoO2, ground state AFM within a Co4+ row ------
\newcommand{\comomentRow}{\ensuremath{0.97}\xspace}       % mu_B on Co4+, ground state
\newcommand{\comomentThreeRow}{\ensuremath{0.00}\xspace}  % mu_B on Co3+, ground state
\newcommand{\comomentRowRange}{\ensuremath{0.97}\ to \ensuremath{1.02}\xspace} % across FM and both AFM
\newcommand{\msatRow}{\ensuremath{0.48}\xspace}           % mu_B per Co if every moment aligned
\newcommand{\dSbound}{\ensuremath{0.23}\xspace}           % k_B per Co, bulk reading of Eq. (S34)
\newcommand{\magorderRow}{antiferromagnetic within a Co\ensuremath{^{4+}} row\xspace}
\newcommand{\efmafmRow}{\ensuremath{28.2}\xspace}         % meV/Co, FM above AFM
\newcommand{\enmRow}{\ensuremath{63.7}\xspace}            % meV/Co, NM above AFM
\newcommand{\gapRow}{\ensuremath{0.70}\xspace}            % eV, AFM ground state (0.698)
\newcommand{\gapRowFM}{\ensuremath{0.43}\xspace}          % eV, FM in the same cell (0.428)
\newcommand{\eafmIntra}{\ensuremath{-56.5}\xspace}        % meV/Co4+, E(AFM) - E(FM), intra-row
\newcommand{\eafmInter}{\ensuremath{+41.7}\xspace}        % meV/Co4+, E(AFM) - E(FM), inter-row
\newcommand{\jintraMag}{\ensuremath{56}\xspace}           % meV, magnitude of the dominant coupling
\newcommand{\dnRow}{\ensuremath{0.20}\xspace}             % e, d-occupation split, ground state
\newcommand{\bondRow}{\ensuremath{0.001}\xspace}          % Angstrom, contrast of site-mean Co-O
\newcommand{\bondSpanRow}{\ensuremath{0.021}\xspace}      % Angstrom, span of individual Co-O bonds

% --- (b) three periodic Li motifs at x = 1/2, 4x2 cell, 8 f.u. ------------
%     Fixed ideal geometry, ferromagnetic alignment imposed.
\newcommand{\deltaEZigzag}{\ensuremath{-8.3}\xspace}       % meV/f.u. vs row, U=4
\newcommand{\deltaEZigzagUzero}{\ensuremath{+8.2}\xspace}  % meV/f.u. vs row, U=0
\newcommand{\deltaEClumped}{\ensuremath{58.1}\xspace}      % meV/f.u. vs row, U=4
\newcommand{\deltaEClumpedUzero}{\ensuremath{92.4}\xspace} % meV/f.u. vs row, U=0
\newcommand{\gapRowMotif}{\ensuremath{0.679}\xspace}       % eV
\newcommand{\gapZigzag}{\ensuremath{0.800}\xspace}         % eV
\newcommand{\gapClumped}{\ensuremath{0.675}\xspace}        % eV
\newcommand{\gapMotifSpread}{\ensuremath{0.13}\xspace}     % eV, 0.800 - 0.675
\newcommand{\dosRowUzero}{\ensuremath{22.5}\xspace}        % states/eV/cell, U=0
\newcommand{\dosContrast}{\ensuremath{1.004}\xspace}       % clumped/row at U=0
\newcommand{\momentRowUzeroA}{\ensuremath{0.3990}\xspace}  % mu_B, row cell at U=0, site 1
\newcommand{\momentRowUzeroB}{\ensuremath{0.3986}\xspace}  % mu_B, row cell at U=0, site 2

% --- (c) sqrt3 x sqrt3 Li_{1/3}CoO2, ground state ferrimagnetic -----------
\newcommand{\comomentXthird}{\ensuremath{1.04}\xspace}     % mu_B on each Co4+
\newcommand{\magorderXthird}{ferrimagnetic, the two Co\ensuremath{^{4+}} moments antiparallel\xspace}
\newcommand{\efmafmXthird}{\ensuremath{12.4}\xspace}       % meV/Co, FM above ferri
\newcommand{\enmXthird}{\ensuremath{230}\xspace}           % meV/Co, NM above ferri
\newcommand{\splitXthird}{\ensuremath{37}\xspace}          % meV per cell, FM above ferri
\newcommand{\gapXthird}{\ensuremath{1.21}\xspace}          % eV, ferri ground state (1.206)
\newcommand{\gapXthirdUfour}{\ensuremath{1.179}\xspace}    % eV, FM at U=4
\newcommand{\gapXthirdUfive}{\ensuremath{1.166}\xspace}    % eV, FM at U=5.5
\newcommand{\dosXthirdNM}{\ensuremath{7.7}\xspace}         % states/eV/cell, NM
\newcommand{\dnXthird}{\ensuremath{0.26}\xspace}           % e, d-occupation split
\newcommand{\bondXthird}{\ensuremath{0.003}\xspace}        % Angstrom, contrast of site-mean Co-O

% --- methodological numbers quoted in the Supplemental Material ------------
\newcommand{\comomentHighU}{\ensuremath{2.28}\xspace}      % mu_B, Co4+ at U=5.5, high spin
\newcommand{\momentFMa}{\ensuremath{1.028}\xspace}         % mu_B, inequivalent FM Co4+, supercell
\newcommand{\momentFMb}{\ensuremath{0.952}\xspace}         % mu_B, the other one
\newcommand{\forcePlateau}{\ensuremath{78}\xspace}         % meV/Angstrom, residual force plateau
\newcommand{\primSign}{\ensuremath{3.1}\xspace}            % meV/Co, primitive-cell FM above NM


% ---------------------------------------------------------------------------
% Figure-or-placeholder helper.  The figure set is produced by a separate pass;
% until a file appears, a labelled box holds its place and the document still
% compiles.  #1 file base name, #2 width.
% ---------------------------------------------------------------------------
\newcommand{\figorpending}[2]{%
  \IfFileExists{figures/#1.pdf}{\includegraphics[width=#2]{#1}}%
  {\fbox{\parbox[c][0.40\linewidth][c]{0.94#2}{\centering\ttfamily\small
   [figure \detokenize{#1}.pdf pending]}}}%
}

% ---------------------------------------------------------------------------
% Float placement.  Two of the column figures are tall and carry long captions,
% so at the LaTeX defaults figure-plus-caption exceeds \topfraction, no top slot
% is ever legal, and the floats queue up and cascade past the bibliography.
% Relaxing the fractions and raising the per-page counters restores top
% placement on the page where each float is discussed.
% ---------------------------------------------------------------------------
\renewcommand{\topfraction}{0.9}
\renewcommand{\bottomfraction}{0.9}
\renewcommand{\textfraction}{0.07}
\renewcommand{\floatpagefraction}{0.75}
\renewcommand{\dbltopfraction}{0.9}
\renewcommand{\dblfloatpagefraction}{0.75}
\setcounter{topnumber}{3}
\setcounter{bottomnumber}{2}
\setcounter{totalnumber}{4}
\setcounter{dbltopnumber}{3}

\newcommand{\kB}{k_{\mathrm{B}}}
\newcommand{\muBs}{\mu_{\mathrm{B}}}
\newcommand{\Ea}{E_{\mathrm{a}}}
\newcommand{\Ton}{T_{\mathrm{on}}}
\newcommand{\Tcol}{T_{\mathrm{ord}}}

\begin{document}

\title{Cold Crystallization of the Lithium Sublattice and the Warming-Only
Resistance Anomaly of Delithiated Li$_x$CoO$_2$}

\author{Santosh Kumar Radha}
\affiliation{Agentfield AI}

\date{\today}

\begin{abstract}
Chemically exfoliated Li$_x$CoO$_2$ nanoflakes carry a resistance anomaly that appears only on
the warming branch of a thermal loop, leaving the noise near 146~K and passing its half height
near 152~K, collapsing through 240 to 245~K in the raw difference between the branches and some
ten kelvin higher once that difference is normalized to a cooling baseline that is itself
falling, gone in either reading by 260~K, and roughly
doubling in a 5~T field while the cooling baseline moves by a few percent. We read the loop as
cold crystallization of the lithium sublattice: nucleation and growth are never favorable
together along a cooling sweep, so the flake reaches base temperature as a disordered matrix
carrying quenched, defect-templated embryos, and those embryos grow as soon as lithium
unfreezes, converting the channel to a more resistive ordered phase that a first-order boundary
later removes. The unfreezing condition fixes an effective migration barrier of
$0.40\pm0.05$~eV, which is the divacancy value in this material, and requires the onset to move
by $12\pm1$~K per decade of warming rate at fixed cooling protocol while the collapse stays
fixed to better than $10^{-3}$~K. Density-functional calculations place charge order, a
low-spin Co$^{4+}$ moment of \comomentRow~$\muBs$ and a hard gap on one transition, and leave
frustration of that charge order by aperiodic lithium as the surviving account of the resistive
contrast. The field-shifted collapse would require the ordered phase to carry the larger
magnetization in a unit of order ten correlated moments, but at $1.4\pm1.4$~K/T the shift is
consistent with zero at one standard deviation, and our own digitization returns the same sign
at $1.0\pm0.5$~K/T in the raw branch difference and at $0.4\pm0.4$~K/T once that difference is
normalized, so that chain and the surface electron gas it
points to remain conditional. A warming-rate series at fixed cooling protocol decides the
central claim with no new instrumentation.
\end{abstract}

\maketitle

\section{Introduction}
\label{sec:intro}

A resistance that depends on which way the thermometer is moving cannot be a property of an
equilibrium state. Chemically exfoliated Li$_x$CoO$_2$ nanoflakes prepared by slow
delithiation show exactly that behavior: the cooling trace from room temperature is smooth and
featureless, while the warming trace of the same device develops an excess resistance near
146~K, reaches a maximum comparable in size to the cooling baseline itself, and collapses
through 240 to 245~K, the raw difference between the two branches passing half its peak height
at 244~K and reaching nothing by 260~K~\cite{Crowley2020,RadhaPerspective2021}. Those collapse
temperatures belong to the raw difference; the same excess normalized to the cooling baseline,
which is the quantity Fig.~\ref{fig:kinetics} plots, falls through its half height about ten
kelvin higher, near 254~K, because the baseline it is divided by is itself decreasing with
temperature, and every collapse number quoted below names the quantity it belongs
to. The loop repeats over thermal cycles
without degradation, so no irreversible chemistry is involved. A perpendicular field of 5~T
roughly doubles the excess while leaving the cooling baseline within a few percent of its
zero-field value, and it pushes the high-temperature collapse upward without moving the
low-temperature onset. Those temperatures, the amplitude and the field response are read from
the published figures of Refs.~\cite{Crowley2020,RadhaPerspective2021} by the digitization
described with Fig.~\ref{fig:kinetics}; 146~K is the foot of the rise and not its midpoint, the
digitized excess crossing a tenth of its peak height at 149~K and half of it at 152~K, and
unlike the collapse the rise gives the same midpoint whichever of the two excesses is used, so
that since the onset condition of Sec.~\ref{sec:kinetics} is a statement about the
half-transformed sample it is 152~K and not 146~K that the model is evaluated at throughout, a
distinction worth about 0.015~eV in the barrier extracted below. The warming rate was not
recorded, so it is carried throughout as an explicit unknown. Reanalysis of the raw traces would
tighten every number quoted below.

Two of those observations already dictate the shape of any explanation. A feature that exists
on one branch of a loop and not the other belongs to a transformation that the cooling sweep
outruns, which in the language of calorimetry is cold crystallization: the transformation is
suppressed on cooling by slow kinetics, is stored in the quenched sample as a metastable
state, and completes on reheating as soon as the relevant degree of freedom regains
mobility~\cite{Wunderlich1958,Wellen2014,Angell1995,Debenedetti2001}. The term belongs to
molecular glasses and polymers~\cite{Wunderlich1958}, to metallic glasses~\cite{Greer1995} and to
the chalcogenide phase-change materials whose crystallization on warming is itself read as a
resistance step~\cite{Wuttig2007,Friedrich2000,Klym2015}, and is used here by analogy, no cation
or vacancy sublattice of an inorganic layered oxide having been described that way before. A
field that acts on the excess and not on the
background belongs to a thermodynamic phase boundary rather than to the carriers of the parent
metal, because a magnetoresistance that doubled the anomaly by acting on the conduction channel
as a whole would have moved the baseline by the same factor; the inference is sound against a
bulk-carrier magnetoresistance and not against one localized at the contacts, which
Sec.~\ref{sec:constraints} takes up separately.

The slow degree of freedom is not hard to identify. Lithium in layered Li$_x$CoO$_2$ migrates
at intermediate concentration by a divacancy hop whose barrier first-principles calculations
place at 0.41~eV, with a scatter of about $\pm0.04$~eV across local lithium and vacancy
environments; the isolated-vacancy path is roughly twice as costly, 0.83~eV, and operates only
near full lithiation~\cite{VanderVen2000,VanderVen2001,VanderVen2001b}. Lithium and vacancy
ordering at $x\simeq1/2$ is established by
diffraction~\cite{Reimers1992,ShaoHorn2003,VanderVen1998}, while the ordering that cluster
expansion predicts at $x\simeq1/3$ has never been observed, its computed boundary lying at or
below room temperature~\cite{VanderVen1998,Wolverton1998,Motohashi2009,Takahashi2007}. Bulk
Li$_x$CoO$_2$ over $0.46\le x\le0.71$ shows a first-order anomaly between 155 and 185~K, placed
near 155~K and nearly independent of composition in one series of crystals and at 175 to 185~K in
another~\cite{Miyoshi2010,Motohashi2009}, together with a distinct and lower lithium-freezing
scale near 130~K below which the magnetization depends on how fast the crystal was cooled and
above which $^{7}$Li motional narrowing and muon-detected diffusion set lithium
moving~\cite{Miyoshi2010,Nakamura2006,Sugiyama2009}; that pair of scales is the bulk counterpart
of the two edges analyzed here. In the sister compound Na$_x$CoO$_2$ the alkali pattern controls
the charge order of the CoO$_2$ sheets directly~\cite{Foo2004,Roger2007}, so an ionic ordering
transition is expected to carry an electronic signature.

This paper builds the minimal model that respects all of those constraints and extracts its
parameters. Section~\ref{sec:constraints} states what the measurement establishes and what the
two-probe geometry prevents it from establishing. Section~\ref{sec:mixture} writes the channel
as a two-phase mixture and fixes the size of the transformation. Section~\ref{sec:kinetics}
derives the unfreezing condition, extracts the barrier, and shows why the cooling branch is
featureless. Section~\ref{sec:field} treats the field as a thermodynamic variable, converts
the observed shift of the collapse into a magnetization difference, and counts how many
correlated moments the field must be addressing. Section~\ref{sec:dft} reports the
density-functional test of the microscopic ingredients, and Sec.~\ref{sec:tests} lists the
measurements that decide the question. An earlier account of this material matched the onset
temperature to an equilibrium electronic crossover computed for a nine-site cluster, and we
disagree with it for a reason internal to the present model, that an onset set by ionic kinetics
moves with sweep rate while no equilibrium electronic scale does.

Throughout, the complete derivations are given in the Supplemental Material~\cite{SM}, whose
equations are numbered S1 onward; the main text carries the physical content and the results.

\section{What the measurement constrains}
\label{sec:constraints}

The device is a chemically exfoliated single-crystal flake of thickness in the 20 to 60~nm
class, contacted by two nickel electrodes across a channel of order 5~$\mu$m on
Si/SiO$_2$~\cite{Crowley2020}, made and read out by the fabrication and transport pipeline the
same laboratory had established on exfoliated oxide nanoflakes~\cite{Crowley2018}. Every contact
metal tried on this material is Schottky-like,
with a quasilinear current-voltage branch appearing only above about 2.5~V, and the resistance
is read from that branch. The measured quantity is therefore
\begin{equation}
  R_{\mathrm{meas}}(T;\mathcal{H})
  = R_{\mathrm{c}}(T,V;\mathcal{H}) + R_{\mathrm{ch}}(T;\mathcal{H}),
  \label{eq:circuit}
\end{equation}
where $R_{\mathrm{c}}$ is the series resistance of the two junctions on that branch,
$R_{\mathrm{ch}}$ is the channel, and $\mathcal{H}$ denotes the thermal history. The excess
\begin{equation*}
  \Delta R(T) \equiv R_{\mathrm{meas}}^{\,\mathrm{warm}}(T)
  - R_{\mathrm{meas}}^{\,\mathrm{cool}}(T)
\end{equation*}
cancels every history-independent term in Eq.~\eqref{eq:circuit}, including any smooth
$R_{\mathrm{c}}(T)$, but it does not cancel a history-dependent contact: lithium rearranging
under the electrode obeys the same kinetics with the same barrier as the channel mechanism, and
would therefore reproduce the branch asymmetry, the amplitude and the rate shift alike. Hysteresis
in a layered-crystal device need not be ionic at all, interfacial charge transfer alone having
been shown to produce a gate-induced loop in van der Waals heterostructures~\cite{Kumar2020}, so
neither the presence of a loop nor its size identifies the mechanism on its own. Two
observations are consistent with a channel origin without settling it. The anomaly correlates
with the delithiation protocol rather than with contact processing, and it survives a change of
contact material, a graphene-contacted device from the same batch showing an anomaly above
1~M$\Omega$ with the same phenomenology~\cite{Crowley2020,RadhaPerspective2021}. Neither
observation carries much weight on its own: that device is a single one, its two edges sit
nearer 185 and 200~K than the 152 and 240~K quoted here, and a larger anomaly obtained with a
different contact metal reads to a skeptic as an effect that scales with contact quality, so
the second observation cuts both ways. Only channel-length scaling and a four-probe measurement
discriminate, neither has been done, and both are listed in Sec.~\ref{sec:tests}. Joule heating
is negligible, since microampere currents at volt-scale bias deposit microwatts into a
substrate-mounted flake.

The channel is also a lithium-ion conductor read at 2.5~V across 5~$\mu$m, so bias-driven
lithium redistribution has to be considered, and slowness is not what excludes it. The
Nernst--Einstein relation applied to the same hop model that fixes the barrier below gives a
diffusivity of $2\times10^{-10}$~cm$^{2}$/s at room temperature, inside the measured chemical
diffusivity band of this material, and a drift velocity of 0.4~$\mu$m/s at the naive channel
field of 5~kV/cm, so an unscreened field would carry lithium millimeters over a single sweep
from 80 to 300~K, a thousand channel lengths. What excludes it is the boundary condition rather
than the rate. Nickel electrodes are ion blocking, so lithium moves only until its own space
charge cancels the applied field, which it does over a Debye length of 0.04~nm within
microseconds above 240~K and at a total displacement of $3\times10^{-14}$~m; the interior of the
channel carries no sustained field and no sustained ionic current. Two Schottky contacts in
series reinforce the conclusion, since most of the applied 2.5~V drops across depletion regions
tens of nanometers wide rather than across the channel. What is not excluded is a Debye-thin
lithium space-charge layer at each electrode modulating its barrier, which is the
history-dependent contact term of Eq.~\eqref{eq:circuit} once more, and a bias-history series,
the amplitude measured against read voltage and against dwell time at bias, is the direct test.
Every quantitative statement below is a statement about $\Delta R$ and never about an absolute
resistivity, which is the limit of what a two-probe experiment establishes.

The anomaly is also selective in preparation. Flakes delithiated slowly, over a three-day soak
in dilute acid, to $x$ between about 0.37 and 0.48 show it; flakes taken rapidly in stronger
acid to nearly the same composition do not, and the onset moves upward as $x$
decreases~\cite{RadhaPerspective2021}. Composition alone therefore does not control the anomaly.
What distinguishes the preparations is how much time the lithium sublattice spent near
equilibrium, which is the first hint that the physics is kinetic. The drift of the onset with
composition is not in conflict with reading it as a clock, because composition reaches the onset
only through the barrier: the divacancy barrier rises as lithium is removed below half filling,
and at the sensitivity $d\Ton/d\Ea=397$~K/eV derived in Sec.~\ref{sec:kinetics} the $\pm0.04$~eV
environment scatter of that barrier~\cite{VanderVen2001b} moves the onset by at most about
$\pm16$~K across the window in play, with the same spread necessarily widening the rise.
The batch series of Table~S2 in Ref.~\cite{SM} turns that into a test.

\subsection{The lithium sublattice in these flakes}
\label{sec:lisublattice}

Nothing below requires the parent phase to be structurally featureless, but everything below
requires its lithium sublattice to be positionally disordered between 250 and 300~K, and that is
what the literature reports at these compositions. The row-ordered monoclinic structure occupies
a narrow field around $x=1/2$, bounded by two-phase plateaus near $x=0.48$ and
$x=0.55$~\cite{Motohashi2009}; room-temperature single-crystal diffraction returns the
rhombohedral solid solution at $x=0.48$ and at $x=0.35$, with the monoclinic distortion
explicitly absent at the first of these and only diffuse superlattice intensity too weak to
refine~\cite{Takahashi2007}; potentiostatically deintercalated single crystals carry no
lithium-vacancy superlattice at the electronically stable composition
$x=0.72$~\cite{OuYang2012}; and polycrystals at $x=0.35$ and
$x=0.40$ are rhombohedral~\cite{Motohashi2009}. The margin is thin even at exact half filling,
where the measured order-disorder temperature is only about 60~$^{\circ}$C, low for an oxide
because charge transfer from lithium to the host screens the lithium-lithium
interaction~\cite{Reimers1992,VanderVen1998}. Four atomic percent of aluminium on the cobalt site
removes the hexagonal-to-monoclinic transition altogether~\cite{ErRami2022}, and an
aluminium-substituted shell reaching only a few hundred angstroms in from the particle surface
does the same~\cite{Cho2001}, so chemical and site disorder is established as a destroyer of this
order while the separate question of finite size is not, and the exfoliation route used here is
known to leave the flakes cobalt deficient~\cite{Pachuta2019}.

The flakes' own material says the same thing directly. Electron diffraction of as-exfoliated
Li$_{0.37}$CoO$_2$ from this preparation finds only the rhombohedral pattern and its slightly
sheared monoclinic companion, the lithium-vacancy superstructure appearing only on annealing
near 200~$^{\circ}$C~\cite{Volkova2021,RadhaPerspective2021}; that anneal runs on a clock
incommensurable with a 1~K/min ramp and is not a measurement of the collapse temperature, but it
fixes both the initial state and the direction of the transformation. Independent chemical
delithiation by the same acid route to $x=0.34$ finds no monoclinic
distortion~\cite{Basch2014}, and restacking single layers destroys interlayer registry by
construction~\cite{Cheng2016}. What forms on warming is therefore not the clean row-ordered
ground state but a partially ordered, domain-textured version of it, which is what the
sub-$\kB\ln2$ latent entropy of Sec.~\ref{sec:field} requires, what the mazed microstructure of
Ref.~\cite{ShaoHorn2003} looks like in real space, and what the coexisting ordered and
disordered regions imaged in Ref.~\cite{Iwaya2013} show.

The sublattice is also mobile in the window where the parent phase has to be disordered.
$^{7}$Li motional narrowing and muon-spin spectroscopy place the onset of lithium hopping between
130 and 150~K~\cite{Miyoshi2010,Nakamura2006,Sugiyama2009}, and the one series that covers a
range of compositions finds it nearly independent of them over $0.46\le
x\le0.71$~\cite{Miyoshi2010}, so by 250~K the ions are moving freely and whatever order they hold
is equilibrated rather than frozen. The sister compound makes the point quantitatively and from
the other side. In Na$_{0.8}$CoO$_2$ the alkali sublattice orders isothermally below an arrest
temperature, with a measured ordering time constant near 40~min at 200~K that lengthens more than
tenfold across the 190 to 210~K interval and a quench below 212~K that suppresses the ordering
entirely~\cite{Schulze2008}. At a different composition and five orders of magnitude higher in
cooling rate, the vacancy liquid that Na$_x$CoO$_2$ carries above 335~K can be frozen on a
nanocalorimeter into what its authors call a sodium vacancy glass~\cite{Galeski2016}, so the
alkali sublattice of this family arrests over a range of protocols rather than at a single
temperature; and the family's high-temperature alkali order-disorder transition, at 427~K in
Na$_{0.66}$CoO$_2$, is sluggish enough to have been followed at more than one heating and cooling
rate~\cite{Ying2008}. Lithium orders more weakly than sodium for the screening reason given
above, which is why sodium cobaltates show well-defined vacancy superlattices where the lithium
compound shows none~\cite{Motohashi2009,OuYang2012}.

One alternative slow degree of freedom deserves disposing of explicitly. Delithiated
Li$_x$CoO$_2$ undergoes a stacking transformation from the O3 sequence to H1-3 and then to O1,
and that transformation is first order, hysteretic and sluggish, so it would produce the same
Kissinger phenomenology as an ionic ordering. It is not available here. The H1-3 phase is
measured near $x=0.12$ and predicted over $0.12$ to $0.19$, O1 appears only at $x=0$, and the O3
host is stable throughout $x>0.3$~\cite{Motohashi2009,VanderVen1998,Chen2002,Amatucci1996}, so
flakes at $x=0.37$ to 0.48 sit a factor of three away from the nearest composition at which a
slab glide is thermodynamically available, and the glide barriers are small only once the
lithium planes are empty. Were one available it would be unmistakable, collapsing the interslab
spacing from 4.81 to 4.52~\AA~\cite{Motohashi2009} and doubling the $c$ axis, against the tenth
of a percent of metric change that lithium ordering produces; that contrast is the content of the two
diffraction entries of Table~S2 in Ref.~\cite{SM}.

\section{The channel as a two-phase mixture}
\label{sec:mixture}

Take the channel to contain two phases: a parent phase P, positionally disordered on the
lithium sublattice as Sec.~\ref{sec:lisublattice} documents, with resistivity
$\rho_{\mathrm{P}}$, and an ordered phase O in which
lithium and vacancy order is accompanied by Co$^{3+}$/Co$^{4+}$ charge order, with resistivity
$\rho_{\mathrm{O}}>\rho_{\mathrm{P}}$. Let $\phi$ be the volume fraction transformed. For any
microgeometry the mixture conductance lies between the series and parallel bounds on a
two-component composite~\cite{Wiener1912,Milton2002}, and both bounds give
\begin{equation}
  \frac{\Delta R}{R_{\mathrm{P}}} = c\,\phi\,(s-1),
  \qquad s \equiv \frac{\rho_{\mathrm{O}}}{\rho_{\mathrm{P}}},
  \label{eq:mixture}
\end{equation}
exactly for the series arrangement with $c=1$, and to first order in $\phi(1-1/s)$ for the
parallel arrangement with $c=1/s$; the derivation and the sense in which $c\in(0,1]$ is a
geometry factor are given in Eqs.~(S1)--(S4) of Ref.~\cite{SM}. Equation~\eqref{eq:mixture}
has the two limits it must have, vanishing when nothing has transformed and vanishing when the
two phases conduct alike.

At the zero-field maximum the excess equals the cooling baseline at the same temperature, so
$c\,\phi(s-1)\simeq1$. The transformation is therefore not a perturbation. Either a volume
fraction of order unity converts with a modest resistivity contrast, or a smaller fraction
converts with a large contrast and is arranged so that the current must cross it. A single
reconstructed surface layer on a flake tens of nanometers thick cannot generically produce a
relative excess of order unity unless that layer already carries at least half the total
conductance, which is precisely what the thickness series of Sec.~\ref{sec:tests} decides. The
volume reading is carried below as the simpler of the two and without prejudice, and the two are
not interchangeable. Equation~\eqref{eq:mixture} itself survives either way, with $\phi$ read as
the transformed fraction of whatever cross-section carries the current, but the per-cobalt
normalizations of Sec.~\ref{sec:field} do not: if the transformation is confined to the
surfaces, the magnetization and entropy densities in Eq.~\eqref{eq:cc} refer to surface cobalt
rather than to all of them, and the thickness dependence of $\Delta R/R_{\mathrm{P}}$ inverts.
Section~\ref{sec:field} keeps both readings open, and Sec.~\ref{sec:tests} gives the pair of
measurements that separates them.

Equation~\eqref{eq:mixture} takes $s>1$ from the measurement and does not explain it, but the
same mixture bounds that produced it also bound how large it can be. In the parallel
arrangement the excess cannot exceed $\phi/(1-\phi)$ however large the contrast grows, so an
excess of order unity already demands that at least half the current-carrying cross-section
transform, and combining that limit with the series one leaves $s$ between about 2 and 5 with
$\phi$ between 0.25 and 0.75. The contrast the data ask for is therefore modest, which by itself
excludes reading the resistivity of the ordered phase as activated across the gap the
calculations of Sec.~\ref{sec:dft} return: half of a 0.7~eV gap would give $s$ between $10^{7}$
and $10^{13}$ across this temperature window, seven to thirteen orders of magnitude from what is
measured, and the published traces decay with an apparent activation scale far below any gap in
the problem~\cite{Crowley2020}. Both phases must therefore conduct through low-energy weight
that a periodic Kohn-Sham gap does not describe, and $s$ between 2 and 5 says that the ordered
phase retains between a fifth and a half of the weight the frozen phase carries. The
first-principles calculations of Sec.~\ref{sec:dft} narrow what can be responsible for even that
much: the contrast is not a property of which periodic lithium arrangement the ordered phase
adopts, since those arrangements are indistinguishable in the one quantity that controls
conduction, and the surviving candidate is the frustration that aperiodic lithium occupation
imposes on the cobalt charge order of the frozen phase. The kinetic and thermodynamic arguments
that follow use only the sign of $s-1$ and are unaffected by which microscopic account of its
size proves correct.

\begin{figure}[!tbp]
  \centering
  \figorpending{fig1_mechanism}{\columnwidth}
  \caption{Schematic of the two-phase free-energy landscape and of the two paths a thermal loop
  takes through it. The functional is drawn for legibility rather than computed, and no
  temperature read off it is a result of this work. (a) Free energy per cobalt as a function of
  the ordering coordinate $\phi$ at three temperatures spanning the first-order boundary
  $\Tcol=240$~K, on an expanded free-energy scale so that the barrier is visible. Above $\Tcol$
  the ordered minimum survives as a metastable well up to a spinodal set by the drawn
  functional, while the disordered parent phase P at $\phi=0$ becomes the global minimum; below
  $\Tcol$ the ordered phase O at large $\phi$ is the deeper of the two. They are separated by
  the nucleation barrier $W^{*}=\pi\lambda^{2}/|\Delta f_{A}|$ of Eq.~(S21), which diverges as
  $T\to\Tcol$ from below and falls as the undercooling grows. (b) Transformed fraction along a
  complete thermal loop. The cooling branch stays at $\phi=0$ because the region in which a
  classical nucleus can form and the region in which it can grow do not overlap along the sweep,
  so the flake reaches base temperature untransformed and carrying the quenched, defect-templated
  embryos of Sec.~\ref{sec:asymmetry}. The warming branch rises at the kinetic onset $\Ton$,
  where those embryos grow, holds near complete conversion, and falls at the thermodynamic
  boundary $\Tcol$ where $\phi_{\mathrm{eq}}$ itself collapses. The onset is a clock reading and
  moves with sweep rate; the collapse is a phase boundary and does not.}
  \label{fig:mechanism}
\end{figure}

\section{Kinetics of the transformation}
\label{sec:kinetics}

\subsection{Unfreezing on warming}

Ordering requires lithium to move. Write the minimal relaxation of the transformed fraction
toward its equilibrium value,
\begin{equation}
  \frac{d\phi}{dt} = k(T)\,\bigl[\phi_{\mathrm{eq}}(T,B)-\phi\bigr],
  \qquad
  k(T) = \nu\,e^{-\Ea/\kB T},
  \label{eq:master}
\end{equation}
with $\nu$ an attempt frequency of order a lithium hop-attempt or phonon frequency, $\Ea$ the
effective barrier, and $\phi_{\mathrm{eq}}$ the equilibrium fraction, equal to unity below the
first-order boundary $\Tcol(B)$ and to zero above it. Detailed balance is respected because
forward and reverse processes share the barrier top, the ratio of their rates being carried
entirely by $\phi_{\mathrm{eq}}$. Section~\ref{sec:asymmetry} explains why
Eq.~\eqref{eq:master} with a single effective barrier is legitimate on the warming branch and
not on the cooling branch.

Start the warming sweep in the quenched state, $\phi=0$ with $\phi_{\mathrm{eq}}=1$, and let
$T(t)=T_{0}+rt$. Integrating Eq.~\eqref{eq:master} and evaluating the resulting integral by
Laplace's method at its upper endpoint, which is controlled by the large parameter
$\Ea/\kB T\simeq31$, gives
\begin{equation}
  \phi(T) = 1-e^{-\Theta(T)},
  \qquad
  \Theta(T) \simeq \frac{\nu}{r}\,\frac{\kB T^{2}}{\Ea}\,e^{-\Ea/\kB T},
  \label{eq:theta}
\end{equation}
with the subleading correction of order $2\kB T/\Ea$, a few percent here; the steps are
Eqs.~(S6)--(S11) of Ref.~\cite{SM}. The dimensionless $\Theta$ is the number of successful
rearrangement attempts accumulated along the sweep, and it scales inversely with the warming
rate, as it must: sweeping twice as fast halves the time spent at every temperature.

Defining the onset $\Ton$ by the rise midpoint, $\Theta(\Ton)=\ln2$, and solving for the
barrier gives
\begin{equation}
  \Ea = \kB\Ton\,
  \ln\!\left[\frac{\nu\,\kB\Ton^{2}}{\ln2\;\Ea\,r}\right],
  \label{eq:Ea}
\end{equation}
an implicit relation solved by iteration in two or three steps. With $\Ton=152$~K, which is
where the digitized excess reaches half its peak height and therefore where the condition
$\Theta=\ln2$ applies, an attempt
frequency $\nu=10^{13}$~s$^{-1}$, and a warming rate $r=1$~K/min, which is a typical
cryostat value carried here as an explicit unknown because the rate was not recorded,
Eq.~\eqref{eq:Ea} returns 0.469~eV if the growth velocity is taken to have saturated at the
onset. It has not saturated there. At 152~K the thermodynamic factor $1-e^{\Delta f/\kB T}$ of
Eq.~\eqref{eq:kjmarates} still stands at 0.12, and the transformation on the warming branch is
growth from a fixed areal density $n_{0}$ of seeds rather than the single-exponential relaxation
of Eq.~\eqref{eq:master}, so the condition that fixes the barrier is $\pi n_{0}
R_{\mathrm{w}}^{2}(\Ton)=\ln2$, with $R_{\mathrm{w}}$ the radius a growth front sweeps along the
ramp. Retaining both corrections lowers the barrier to 0.403~eV at $n_{0}=10^{-2}$~nm$^{-2}$,
the two quantities entering together as
$\Ea=0.403+0.015\log_{10}(n_{0}/10^{-2}\,\mathrm{nm}^{-2})$~eV, so the barrier has to be quoted
as one member of a pair and never on its own. Propagating the uncertainties gives
\begin{equation*}
  \Ea = 0.40 \pm 0.03\,(\nu) \mp 0.03\,(r) \pm 0.02\,(\Ton) \pm 0.01\,(n_{0})\ \mathrm{eV},
\end{equation*}
per decade of attempt frequency, per decade of warming rate, for $\pm7$~K on the onset and per
decade of seed density, which combine in quadrature to $\Ea=0.40\pm0.05$~eV with a
corner-to-corner range of 0.28 to 0.48~eV; the accompanying $n_{0}$ runs from $10^{-4}$ to
$10^{-1}$~nm$^{-2}$, which is seeds 3 to 100~nm apart, or $10^{2}$ to $10^{5}$ of them per
square micrometer.

That value lands on a mechanism rather than merely inside a band. The divacancy hop that carries
lithium at intermediate concentration has a computed barrier of 0.41~eV with an
environment-dependent scatter of $\pm0.04$~eV, and an average over the composition
window $x=0.37$ to 0.50 that these flakes occupy which has been read as low as 0.36 to 0.39~eV,
against 0.83~eV for the isolated-vacancy path
that operates only near full lithiation~\cite{VanderVen2001b}; the extracted 0.403~eV sits
between those two readings rather than on either of them, which is as much as should be claimed
for it, since what the onset locates is the divacancy family and not one digit inside it, while
the isolated-vacancy path is excluded by a factor of two on any reading.
The identification should not be pressed
hard. With $\nu$, $r$ and $n_{0}$ all free, the quoted interval very nearly fills the range of
plausible lithium barriers in this material, so the coincidence has little discriminating power
until the warming-rate series of Sec.~\ref{sec:tests} pins the slope. What it does support is
reading the onset as an ionic clock rather than as an electronic energy scale.

\subsection{The rate shift, and why it is parameter-free}

Because $\Ea$ enters Eq.~\eqref{eq:Ea} only inside a logarithm, the onset is insensitive to the
prefactors and sensitive to the sweep rate in a way that is easy to measure. Differentiating
the condition $\Theta(\Ton)=\ln2$ at fixed $\Ea$ and $\nu$ gives
\begin{equation}
  \frac{d\Ton}{d\ln r} = \frac{\kB\Ton^{2}}{\Ea+2\kB\Ton},
  \label{eq:rateshift}
\end{equation}
which is the Kissinger relation for a first-order rate process~\cite{Kissinger1957}. At
$\Ton=152$~K and $\Ea=0.40$~eV this is 4.7~K per $e$-fold. The thermodynamic factor of
Eq.~\eqref{eq:kjmarates} carries its own temperature dependence at the onset,
$d\ln f/dT=0.017$~K$^{-1}$ there, which steepens the slope by 8.3~percent, slightly larger than
the $2\kB T/\Ea$ correction the Laplace evaluation already retains and of the opposite sign; keeping
both gives 11.6~K per decade of warming rate, and across the barrier range admitted above the
prediction spans 10.8 to 12.6~K per decade, which is $12\pm1$~K per decade to the precision the
model supports.

That prediction holds only if the cooling protocol is held fixed while the warming rate is
varied. Should the cooldown be slaved to the warm-up, athermal seeding leaves the slope where
it is, but thermally born seeds would thin as $n_{0}\propto1/|r|$ and steepen the shift to about
17.4~K per decade, so a series run without that precaution measures a number whose
interpretation depends on the very seeding mechanism the series is meant to probe. Subject to the
protocol, the shift is parameter-free in the sense that matters: the same differentiation
applied to a transformation whose exponent is $\Theta^{n}$ rather than $\Theta$ returns
Eq.~\eqref{eq:rateshift} unchanged, because the factor $n$ multiplies the rate dependence and
the temperature dependence alike and cancels, which the algebra of Eq.~(S17) of
Ref.~\cite{SM} shows and which we have confirmed numerically for $n=1$ and $n=2$. The predicted
shift therefore does not depend on whether the transformation is growth-limited from quenched
seeds or governed by continuous nucleation, which makes it the sharpest single test of the whole
picture.

The width of the rise is not equally insensitive. Growth in two dimensions from a fixed density
of seeds is quadratic in $\Theta$, and at this barrier such a transformation spans its 10 to 90
percent range over 7.7~K, from 147.4 to 155.1~K. The measured width, $19^{+6}_{-9}$~K, sits well
above that, and a
spread of barriers of about $\pm0.044$~eV closes the gap at the sensitivity
$d\Ton/d\Ea=397$~K/eV, which compositional inhomogeneity of the kind seen across
batches~\cite{RadhaPerspective2021} would naturally supply. The same spread is what limits how
far the onset can move with composition, so the width of the rise and the batch dependence of
the onset are one statement rather than two. Strongly cooperative kinetics, which would sharpen
the rise, is not required by the data and is mildly disfavored by it.

\subsection{Why the cooling branch is featureless}
\label{sec:asymmetry}

Equation~\eqref{eq:master} is symmetric under reversal of the sweep and cannot by itself
produce an anomaly on one branch only, which is the step most discussions of thermal hysteresis
leave out. Applied to the cooling branch,
Eq.~\eqref{eq:master} would predict that the sample transforms immediately below $\Tcol$, where
the relaxation time $\nu^{-1}e^{\Ea/\kB T}$ is at most $3\times10^{-4}$~s across the admitted
barrier range and shorter than $10^{-4}$~s at its center. The asymmetry
therefore cannot come from a single rate; it comes from the two-step structure of a real
first-order transformation, in which a nucleus must appear before anything can grow.

For two-dimensional ordering in the lithium plane the nucleation rate per unit area and the
growth velocity are
\begin{equation}
\begin{aligned}
  I_{\mathrm{n}} &= I_{0}\,e^{-[\Ea+W^{*}(T)]/\kB T},\\
  v &= v_{0}\,e^{-\Ea/\kB T}\bigl(1-e^{\Delta f/\kB T}\bigr),
\end{aligned}
  \label{eq:kjmarates}
\end{equation}
with $\Delta f=f_{\mathrm{O}}-f_{\mathrm{P}}$ the free-energy difference per cobalt, negative
below $\Tcol$, and $W^{*}(T)=\pi\lambda^{2}/|\Delta f_{A}(T)|$ the barrier to forming a
critical disk of ordered phase, $\lambda$ being the domain-wall line tension and $\Delta f_{A}$
the free-energy difference per unit area. The growth velocity vanishes at $\Tcol$, as it must,
and saturates at $v_{0}e^{-\Ea/\kB T}$ deep in the ordered regime. Feeding
Eq.~\eqref{eq:kjmarates} into the Kolmogorov--Johnson--Mehl--Avrami construction for the
transformed fraction~\cite{Kolmogorov1937,Johnson1939,Avrami1939,Avrami1940}, given as Eqs.~(S23)--(S25)
of Ref.~\cite{SM}, yields the asymmetry directly.

The two factors in Eq.~\eqref{eq:kjmarates} pull in opposite directions, and that opposition is
the whole argument, but the place where they cross has to be located correctly. Nucleation needs
a large undercooling, because $W^{*}$ diverges at $\Tcol$ and falls only as the driving force
grows; growth needs a high temperature, because the mobility factor $e^{-\Ea/\kB T}$ collapses
on cooling. The nucleation rate, however, carries that same mobility factor, so it is not large
at the bottom of the sweep: below about half the collapse temperature $I_{\mathrm{n}}$ increases
monotonically with temperature, and nucleation at base temperature is the least likely event in
the problem rather than the most. The only window on the cooling branch in which a classical
nucleus can form at all is the nose of the transformation curve, which for every admissible line
tension sits between 178 and 211~K. A nucleus born there sweeps between 0.6~nm and 40~$\mu$m
through the remainder of a 1~K/min cooldown, while the warming branch needs its seeds to grow
only 1.7 to 20~nm to reach half conversion at 152~K, so a nucleus born at the nose would
transform the flake on the way down and destroy the very asymmetry it was invoked to explain.
The requirement is quantitative and nearly independent of the barrier: every seed present at
base temperature must have been born below $\Ton-9$~K, which is below about 143~K.

Classical nucleation cannot put them there. Forcing the maximum of $I_{\mathrm{n}}$ below 143~K
demands $W^{*}/\kB T\simeq60$ at that temperature, which suppresses the rate by $e^{-60}$, and
a scan over the barrier from 0.34 to 0.45~eV, the attempt frequency over three decades, the
latent entropy from 0.1 to 0.69~$\kB$, the line tension over four decades and the cooling rate
from 1 to 60~K/min leaves the delivered seed density short of what the warming branch requires
by five to eight decades at 1~K/min, and by 1.9 decades in the single most favorable corner.
Heterogeneous nucleation does not repair it, since a potency factor multiplying $W^{*}$ is the
same one-parameter family as a reduced line tension while its prefactor is smaller by the areal
density of the sites themselves, so the shortfall grows by about a decade rather than closing.
Formal closure would need a cooldown of order $10^{3}$~K/min, which is a plunge quench and not a
cryostat.

What the model needs, and what we adopt as an assumption rather than derive, is that the seeds
are not thermal at all. Any barrier of the classical form $W^{*}\propto1/(\Tcol-T)$ is excluded
by the argument above, so the escape has to break the functional form: a quenched population of
defect-templated embryos, present in the flake before the loop begins, each of which goes
supercritical when the critical radius $R^{*}=\lambda/|\Delta f_{A}|$ falls to its own size,
with no mobility factor to pay. At 143~K that radius is 0.7 to 1.4~nm, twenty to ninety cobalt
sites, and the assumption carries an explicit condition that should be stated with it: the
embryo distribution must be cut off above about 1.4~nm, since one 2~nm embryo anywhere in the
flake would go supercritical near 200~K and sweep ten micrometers. Such a cut-off is natural for
patches templated by the defect population that exfoliation leaves behind, whose sizes are set
by the defects and not by the temperature, and unnatural for thermal fluctuations of the parent
phase. On warming, nucleation has in either case already been paid for, the only remaining
barrier is the one that controls growth, and that is exactly the situation
Eq.~\eqref{eq:master} describes with a single $\Ea$. This is cold crystallization, and it
accounts for the branch asymmetry with no electronic ingredient whatever.
Figure~\ref{fig:mechanism} draws the two paths.

That assumption is the weakest joint in the account, and it converts into two predictions, both
of them cheap. If the seeds are athermal the
onset cannot depend on how fast the flake was cooled, whereas thermally born seeds accumulate as
$n_{0}\propto1/|r_{\mathrm{cool}}|$ and would raise the onset by 5.8~K per decade of cooling
rate, the faster cooldown delivering the fewer seeds and hence the higher onset; one device and
one afternoon settle the provenance of the seeds. And because the area swept on the cooling
branch grows as $|r_{\mathrm{cool}}|^{-3}$, a cooldown slower than about 0.1~K/min must
transform the flake on the way down and erase the warming anomaly altogether, which is the
sharpest consequence of the picture that nobody has yet looked for.

The picture also makes an isothermal prediction. With seeds already present, a hold anywhere on
the warming branch is growth limited and its half-conversion time is the radius the warming
branch needs divided by the growth velocity, so a hold at 141 to 146~K, six to eleven kelvin
below the onset, should complete in between three minutes and two hours, centrally between
sixteen and forty-six minutes, and a hold that produces nothing there is genuinely falsifying. A
hold between 150 and 170~K on the way down is a different measurement, because it probes the
nucleation rate rather than growth, and the model supplies only a lower bound on its timescale,
at least one to ten hours at 1~K/min; a null result there excludes nothing. The sister compound
supplies the empirical anchor the model cannot: in Na$_x$CoO$_2$ the alkali sublattice orders
isothermally with a measured time constant near 40~min at 200~K, and a quench below 212~K
suppresses the ordering entirely~\cite{Schulze2008}.

\subsection{The high-temperature collapse}

Above the formation window the ordered phase is in quasiequilibrium, since its relaxation time
at 240~K is at most $3\times10^{-4}$~s, shorter than any laboratory sweep by seven orders of
magnitude,
so the decay of the anomaly is thermodynamic rather than kinetic. Combining Eq.~\eqref{eq:mixture} with the equilibrium
fraction gives
\begin{equation}
  \Delta R(T) = c\,R_{\mathrm{P}}(T)\,\phi_{\mathrm{eq}}(T,B)\,\bigl[s(T,B)-1\bigr],
  \label{eq:decay}
\end{equation}
in which two factors shrink as the temperature rises: the transport contrast $s-1$ of the more
resistive phase measured against an increasingly conductive parent, and the equilibrium fraction
itself, which collapses at $\Tcol(B)$ over a range broadened by the spread of local
compositions. The published traces~\cite{Crowley2020} decay with an apparent activation scale
far below any gap in the problem, which
is what Eq.~\eqref{eq:decay} leads one to expect, since it convolves three temperature
dependences and none of them is an activated conductivity on its own. No gap should be read from
that slope, and quantifying it calls for the raw traces rather than published figures; still
less should the gap the calculations of Sec.~\ref{sec:dft} return be read into it, since the
modest contrast the mixture bounds permit is incompatible with activation across 0.7~eV by seven
orders of magnitude and more. What is
interpretable without them is the ordering of the two branches: the 5~T decay is steeper and
sits higher in temperature, which is what a field that raises $\Tcol$ and increases the contrast
produces, and which a field acting on the ionic kinetics could not produce, since that would
have moved the onset.

The collapse places $\Tcol(0)=240\pm10$~K for this flake, the raw branch difference passing
through 240 to 245~K and reaching zero by 260~K, with the normalized excess of
Fig.~\ref{fig:kinetics} collapsing about ten kelvin higher for the arithmetic reason that its
denominator falls with temperature as well; the parameter that enters
Eqs.~\eqref{eq:kjmarates} and \eqref{eq:decay} is the boundary itself and is fixed by the raw
difference, which is the quantity that carries no baseline in it. The bulk comparison has to be
made against the right scale, because bulk Li$_x$CoO$_2$ carries two. The lower one, near 130~K, is
kinetic: below it the magnetization depends on how fast the crystal was cooled, and it coincides
with the temperature at which $^{7}$Li motional narrowing and muon-detected diffusion report the
lithium ions coming to rest~\cite{Miyoshi2010,Nakamura2006,Sugiyama2009}. The upper one, between
155 and 185~K, is thermodynamic: it is first order and hysteretic, and the differential scanning
calorimetry that places it at 175 to 185~K returns a latent heat of 82~J/mol at $x=0.50$ and
272~J/mol at $x=0.67$~\cite{Motohashi2009}.
Mapped onto that pair, the flake's kinetic edge at 152~K and the bulk kinetic scale agree with
nothing adjusted, and only the thermodynamic edge is displaced, 240~K against 155 to 185~K. Which
flake edge is kinetic and which is thermodynamic is therefore not in question; what needs
explaining is a single shifted phase boundary.

Four candidates for that shift can be ranked, and one that is often offered should be dropped.
Composition comes first and is the only one already evidenced: the ordered volume fraction in
bulk crystals shrinks continuously below $x\simeq0.45$ and the ordering field itself is bounded
near $x=0.48$~\cite{Miyoshi2010,Motohashi2009}, so flakes at $x=0.37$ to 0.48 do not sit on the
bulk boundary at all but drift toward the $\sqrt3\times\sqrt3$ arrangement at $x=1/3$, whose
energy scale Sec.~\ref{sec:dft} computes separately. Second is the cobalt-site disorder that the
exfoliation route introduces, which is a mechanism and not an analogy: exfoliation removes
cobalt along with lithium~\cite{Pachuta2019}, and four atomic percent of aluminium on the cobalt
site is enough to remove the hexagonal-to-monoclinic transition of the bulk
entirely~\cite{ErRami2022}. Third is surface stabilization of the order, for which there
is independent precedent, two groups having concluded that the 155 to 185~K anomaly of
chemically delithiated material is a surface-stabilized phase of significant thermal hysteresis
and aging character, absent from homogeneous potentiostatic single
crystals~\cite{Hertz2008,OuYang2012}. Proton co-intercalation from the acid route is fourth,
with evidence on both sides, proton insertion at deep delithiation having been measured by
prompt gamma activation analysis~\cite{Choi2006} while hydrogen evolution during acid
delithiation argues against an exchange~\cite{Takahashi2007}, and no proton content has been
measured in these flakes; strain we keep only as a formal possibility, a flake lying unclamped
on SiO$_2$ having none measured. Dimensional confinement, on the other hand, carries the wrong
sign and should not be listed: finite-size scaling depresses an ordering temperature rather than
raising it by 40 percent, so what can help is a specific surface enhancement, which is the
third candidate, and not confinement itself.

The model separates the two edges by their rate dependence rather than by their temperature.
$\Ton$ is fixed by $\Ea$, $\nu$, $r$ and the seed density $n_{0}$, and composition reaches it
only through $\Ea$, so a batch series across $x=0.37$ to 0.48 should move it by at most about
$\pm16$~K and should widen its rise by the same amount, while $\Tcol$ tracks the phase boundary
and may move by tens of kelvin. A batch series therefore moves the two edges by very different
amounts, and a warming-rate series moves only one of them at all.

\begin{figure}[!tbp]
  \centering
  \figorpending{fig2_kinetics}{\columnwidth}
  \caption{Model resistance loops from the kinetic equations, set against the published traces.
  (a) Excess $\Delta R/R_{\mathrm{P}}$ obtained by carrying the growth velocity of
  Eq.~\eqref{eq:kjmarates} through the transformed fraction along a
  cooling and a warming sweep with $\Ea=0.4026$~eV, $\nu=10^{13}$~s$^{-1}$ and $\Tcol=240$~K, and
  converting to resistance through Eqs.~\eqref{eq:mixture} and \eqref{eq:decay}, at three
  warming rates. Overlaid is the measured loop of the $x=0.48$ device of Ref.~\cite{Crowley2020},
  digitized from its Fig.~S5 and expressed in the same variable,
  $(R_{\mathrm{warm}}-R_{\mathrm{cool}})/R_{\mathrm{cool}}$, so that the cooling branch is the
  zero of the plot by construction; nothing is fitted, and the model curve carries only the
  parameters listed above. Because that variable divides by a baseline which is itself falling
  with temperature, its rise midpoint at 152~K is the same as the one the raw difference
  $R_{\mathrm{warm}}-R_{\mathrm{cool}}$ gives, while its collapse sits about ten kelvin above
  the 240 to 245~K through which the raw difference passes; the collapse temperatures quoted in
  the text belong to the raw difference and the curves plotted here to the ratio.
  The 5~T trace of the same figure is drawn for reference and is not
  described by a zero-field model. The model cooling branch, which is common to all three rates
  and is therefore drawn once, stays flat for as long
  as the sweep outruns the nose of Sec.~\ref{sec:asymmetry}, and the bar marks the 10 to 90
  percent rise width of the 1~K/min warming curve, 7.7~K about an onset of 152~K. The model peak
  sits somewhat above the
  $c\phi(s-1)\simeq1$ of Sec.~\ref{sec:mixture} because it is integrated at fixed contrast rather
  than normalized to that condition. (b) Rise midpoints against $\log_{10}r$ over more than two
  decades of warming rate at fixed cooling protocol. The line is Eq.~\eqref{eq:rateshift} with
  the thermodynamic factor retained, giving 11.6~K per decade at $\Ea=0.4026$~eV and containing no
  parameter beyond the barrier already fixed by the onset itself; the collapse temperature, being
  thermodynamic, does not move at all. Separating these two slopes is the measurement that
  distinguishes a kinetic onset from an equilibrium one.}
  \label{fig:kinetics}
\end{figure}

\section{The magnetic field and the ordered phase}
\label{sec:field}

\subsection{What the shifted collapse implies}

Treat the collapse as a first-order line in the temperature-field plane, defined by
$\Delta f(\Tcol,B)=0$. Equating the differentials of the two free energies along that line,
with entropy and magnetization per cobalt, gives the Clausius--Clapeyron relation in its
magnetic form,
\begin{equation}
  \frac{d\Tcol}{dB} = \frac{\Delta M(B)}{\Delta S},
  \qquad
  \begin{aligned}
    \Delta M &\equiv M_{\mathrm{O}}-M_{\mathrm{P}},\\
    \Delta S &\equiv S_{\mathrm{P}}-S_{\mathrm{O}},
  \end{aligned}
  \label{eq:cc}
\end{equation}
derived as Eqs.~(S30)--(S32) of Ref.~\cite{SM}. The latent entropy $\Delta S$ is positive
because the ordered phase is the one with fewer configurations available to it, so an upward
shift of the collapse in field requires $\Delta M>0$: the ordered, more resistive phase would
then be the one carrying the larger magnetization. That is the opposite of what happens in
the manganites, where a field melts charge order because the competing metallic phase is the
ferromagnet~\cite{Kuwahara1995,Tokura2006}; here the field stabilizes the resistive phase, so
the moments have to live inside it.

The measured shift, $\Delta\Tcol/\Delta B=1.4\pm1.4$~K/T, converts with
$\muBs/\kB=0.672$~K/T into
\begin{equation}
  \Delta M = (2.1\pm2.1)\,\frac{\Delta S}{\kB}\;\muBs\ \text{per Co}.
  \label{eq:dM}
\end{equation}
The error bar is dominated by the uncertainty in reading the collapse from published figures,
and at one standard deviation the shift is consistent with zero, so its sign is not yet
established. Our own digitization of the same figure moves that conclusion a little and does not
overturn it. Taking the midpoints of the two falls in the raw branch difference gives a shift of
$5\pm2$~K over 5~T, that is $1.0\pm0.5$~K/T, which agrees with the published reading and reaches
two standard deviations from zero; taking them in the excess normalized to the cooling baseline
gives only $2\pm2$~K, that is $0.4\pm0.4$~K/T, which does not. The sign is positive in both
conventions and the magnitude is smaller than the published value in both, so nothing downstream
needs adjusting, since every bound in this section is derived from the larger reading and is
therefore the tighter one, and the honest summary is that the digitization strengthens the sign
determination in the raw difference without settling it in the ratio.
What the present data do establish robustly is
the other field fact, that 5~T roughly doubles the amplitude of the excess while leaving the
cooling baseline within a few percent of itself, and it is the series at 1, 2, 3, 5, 7 and 9~T
listed in Sec.~\ref{sec:tests} that turns the chain from conditional into quantitative.

Equation~\eqref{eq:dM} constrains a ratio of densities. Both $\Delta M$ and $\Delta S$ are per
cobalt of the phase that transforms, so what the measured slope means depends on where the
transformation lives, and the two possibilities differ by two orders of magnitude. Read as a
bulk transformation, the ceiling is set by the Co$^{4+}$ count: at $x\simeq1/2$ half the cobalt
sites are Co$^{4+}$ in a low-spin $d^{5}$ configuration carrying \comomentRow~$\muBs$ apiece, as
Sec.~\ref{sec:dft} confirms, so the largest magnetization available is \msatRow~$\muBs$ per Co
with every moment aligned, and Eq.~\eqref{eq:dM} bounds the latent entropy at
$\Delta S\lesssim\dSbound\,\kB$ per Co, about a third of the $\kB\ln2$ that complete lithium
order and disorder at half filling would release. That reading carries a difficulty of its own,
which Sec.~\ref{sec:dft} sets out and does not resolve: the computed bulk ground states carry no
net moment, so the ceiling is never approached and a purely bulk account of the field shift is
not available in the form the calculation returns.

Read instead as a transformation whose magnetization resides in the surface electron gas, the
same slope means something different. A flake 20 to 60~nm thick holds forty to a hundred and
twenty CoO$_2$ layers between its two surfaces, so a magnetization of order one Bohr magneton
per surface cobalt is about a hundredth of that per cobalt of the whole sample, and the latent
entropy per bulk cobalt implied by Eq.~\eqref{eq:dM} falls by the same factor. The bulk reading
predicts a latent entropy of order a tenth of $\kB$ per cobalt, which is a latent heat
$T\Delta S=2.1$~meV per cobalt, 0.20~kJ per mole of cobalt, or 2.1~J/g, some twenty times a
routine differential-scanning-calorimetry sensitivity of 0.1~J/g and of the same order as the 82
and 272~J/mol measured on bulk crystals at $x=0.50$ and $x=0.67$~\cite{Motohashi2009}; the
surface reading predicts 0.02~J/g, five times below that sensitivity. Since the two also make
opposite predictions for how the excess scales with thickness, calorimetry and the thickness
series form a discriminating pair rather than two separate checks, and Sec.~\ref{sec:tests}
states them that way.

Neither reading allows the parent metal to be the source. Pauli paramagnetism cannot supply a
magnetization of this size, and the magnetic phase diagram of bulk Li$_x$CoO$_2$ places the
boundary between a Pauli-paramagnetic and a Curie-Weiss metal between $x=0.35$ and
$x=0.40$~\cite{Motohashi2009}, overlapping the composition window in which the anomaly appears.
The transformation is in either case a partial ordering, since the entropy it releases is at
most about a third of the $\kB\ln2$ that complete lithium order and disorder at half filling
would give. Two nearby statements should be kept apart from that one. The broadened rise and the
compositional spread of the batches are a width and not an entropy, and the near-degenerate
manifold of stripe variants found in Sec.~\ref{sec:dft} is sub-extensive, accounting for the
domain texture of the ordered phase and for no part of the entropy deficit.

\subsection{How many moments the field addresses}

The magnetization inferred from Eq.~\eqref{eq:dM} cannot be delivered by independent spins. At
5~T the Zeeman energy of a single $g=2$ moment is $g\muBs B=0.579$~meV, while the thermal energy
over the window in which the field effect operates runs from $\kB T=12.9$~meV at 150~K to
20.7~meV at 240~K, so a single spin is outmatched by a factor between 22 and 36 and is polarized
to about two percent. What follows from that is a question of polarization rather than of energy
balance, since Eq.~\eqref{eq:cc} has already exhausted the free-energy comparison. A block of
$N$ moments locked parallel by exchange polarizes as $m=\rho\tanh(N\muBs B/\kB T)$, with
$\rho=\msatRow~\muBs$ per cobalt the moment density that Sec.~\ref{sec:dft} computes, and
requiring it to reach the magnetization $2.1\,\sigma$ that Eq.~\eqref{eq:dM} demands of a latent
entropy $\sigma=\Delta S/\kB$ per cobalt gives
\begin{equation}
  N \simeq \frac{\kB T}{\muBs B}\,
  \operatorname{arctanh}\!\left(\frac{2.1\,\sigma}{\rho}\right),
  \label{eq:headcount}
\end{equation}
which at 5~T and 150~K, where the ordered phase forms, returns $N\simeq10$ for $\sigma=0.05$ and
$N\simeq21$ for $\sigma=0.10$, rising to 16 and 34 at the 240~K collapse. The number to carry
forward is that the field-coupled object holds of order ten aligned moments rather than one,
with a range of five to twenty-five. The ceiling $\sigma\lesssim\dSbound$ used above is
not an independent fact about the material but the $N\to\infty$ limit of
Eq.~\eqref{eq:headcount}, the point at which the block is fully polarized and further moments
buy nothing, so the two statements are two ends of one curve rather than two findings that
corroborate each other. Two structures in this material could supply such a block. One is a
ferromagnetically correlated Co$^{4+}$ domain inside the charge-ordered phase, which
Eq.~\eqref{eq:headcount} sizes at a few lattice constants across. The other is the
spin-polarized two-dimensional electron gas predicted to form on CoO$_2$-terminated surfaces of
this material~\cite{Radha2021}, on whose surface angle-resolved photoemission independently finds
an electronic reconstruction of the kind a polar termination
demands~\cite{Okamoto2017}, a native reservoir of polarizable moments in a band whose
exchange splitting is already built in, load-bearing here for its spin physics rather than for
the hopping scale it also supplies. Only the second survives Sec.~\ref{sec:dft}, which finds the
bulk Co$^{4+}$ moments held antiparallel by an exchange a hundred times the 5~T Zeeman energy.

\subsection{Which coupling, and how to tell}

The field roughly doubles the anomaly while moving the cooling baseline by a few percent and
leaving the onset where it was. Read through Eq.~\eqref{eq:decay}, with $\phi$ on the rising
side field-independent and $R_{\mathrm{P}}$ nearly so, the doubling requires the transport
contrast of the ordered phase itself to double,
\begin{equation}
  \frac{s(T,5\,\mathrm{T})-1}{s(T,0)-1} \simeq 1.9\ \text{to}\ 2.1 ,
  \label{eq:contrast}
\end{equation}
over and above the upward shift of $\Tcol$ already counted. Doubling an activated conductivity
ratio asks for an extra $\kB T\ln2$ of transport gap, which is 9~meV at 150~K and 11~meV at
183~K, a laboratory-scale energy and not an implausible one for a field acting on a correlated
domain of order ten moments.

Four mechanisms can produce it, and they are distinguished by their field dependence and
anisotropy rather than by their size, as Fig.~\ref{fig:field} sets out. A field that enhances the
order parameter or aligns
domains through their net moment gives a Brillouin-like growth of the excess that saturates,
together with a nearly isotropic response. A field that stabilizes the ordered phase through an
induced moment contributes $-\tfrac12\Delta\chi B^{2}$ to the free-energy difference and
therefore grows quadratically without saturating at laboratory fields. A field acting on a
parent network near its percolation threshold~\cite{Kirkpatrick1973} amplifies a small change
in connectivity nonlinearly, and predicts threshold behavior, sensitivity to the sweep
protocol, and large device-to-device scatter. Any mechanism tied to flux through carrier
orbits carries a $\cos\theta$ dependence on field angle, which the spin mechanisms do not. Only
perpendicular-field data exist, so the angular test has never been performed. A field series at
1, 2, 3, 5, 7, and 9~T under a fixed thermal protocol, together with one in-plane run,
separates all four with no adjustable parameters, and is the single most informative
measurement available on this system.

\begin{figure*}[!t]
  \centering
  \figorpending{fig3_field}{\textwidth}
  \caption{Field thermodynamics of the ordered phase and the measurements that identify the
  coupling. (a) Clausius--Clapeyron construction of Eq.~\eqref{eq:cc}. The collapse line
  $\Tcol(B)$ tilts upward with slope $\Delta M/\Delta S$, separating the ordered phase O, which
  carries the larger magnetization, from the parent phase P; the contours are lines of constant
  $\Delta M$ per cobalt at a latent entropy of $0.2\,\kB$ per cobalt, near the ceiling
  Eq.~(S35) permits, and the measured shift, whose one-standard-deviation interval includes
  zero, selects a band of them rather than one.
  (b) The polarization construction of Eq.~\eqref{eq:headcount}. The magnetization
  $\rho\tanh(N\muBs B/\kB T)$ of a block of $N$ exchange-locked moments at 5~T and 150~K is drawn
  against $N$, with the moment density $\rho=\msatRow~\muBs$ per cobalt as its ceiling; the
  horizontal lines are the magnetizations $2.1\,\sigma$ that Eq.~\eqref{eq:dM} requires at latent
  entropies of $0.05$ and $0.10\,\kB$ per cobalt, and they are crossed at $N\simeq10$ and
  $N\simeq21$; the dashed curve repeats the construction at 183~K, the upper end of the
  anomaly, where thermal disordering pushes the same two crossings out to $N\simeq12$ and
  $N\simeq26$. The ceiling is approached only as $N\to\infty$, which is the $\dSbound\,\kB$ bound
  of Eq.~(S35) seen from the other end rather than an independent statement. (c) The predicted
  shapes of the excess amplitude against field, all normalized to agree at the one field that has
  been measured so that only the shape distinguishes them: Brillouin curves for aligned permanent
  moments at $N=20$ and $N=100$, a quadratic curve for an induced moment, and a threshold curve
  for spin-dependent percolation. Saturation is weak at the smaller block size, the excess at 9~T
  exceeding that at 5~T by a factor 1.64 against the linear value 1.80, while $N=100$ gives 1.05.
  The shaded band is the precision the proposed series would have to reach: at $\pm4.5$~percent,
  which is $\pm3.2$~percent on each measured amplitude, all four cases separate at two standard
  deviations, whereas at a more realistic $\pm10$~percent the series still separates $N=20$ from
  $N=100$ and from the quadratic case but not from the linear limit, so resolving weak saturation
  is the demanding part of the measurement.}
  \label{fig:field}
\end{figure*}

\section{Microscopic structure from first principles}
\label{sec:dft}

The model of Secs.~\ref{sec:mixture} to \ref{sec:field} rests on three microscopic statements,
each a well-posed density-functional question. There must be a lithium ordering transition at
all, with an energy scale that puts its boundary near the observed collapse. The ordered phase
must carry Co$^{4+}$ local moments, since Eq.~\eqref{eq:dM} asks it for a magnetization and
Eq.~\eqref{eq:headcount} asks that the magnetization be collective. And it must be the more
resistive of the two phases, which is the one point where Eq.~\eqref{eq:mixture} touches the
electronic structure. The first two statements come back confirmed, the third does not, and the
way in which it fails changes what the transport model is allowed to claim.

The calculations are spin-polarized, in the generalized gradient
approximation~\cite{PBE1996} with projector augmented-wave datasets~\cite{Blochl1994} and a
Hubbard $U$ on the cobalt $3d$
shell~\cite{Anisimov1991,Dudarev1998,Cococcioni2005,Zhou2004,Wang2006}, as implemented in
\textsc{Quantum ESPRESSO}~\cite{Giannozzi2009,Giannozzi2017}, with Marzari--Vanderbilt
smearing~\cite{Marzari1999} and plane-wave cutoffs of 60 and 480~Ry. Three cells are treated:
the row-ordered $x=1/2$ structure of Reimers and Dahn~\cite{Reimers1992,VanderVen1998} with its
antiferromagnetic supercells; a $4\times2$ supercell of eight formula units in which the row
arrangement is compared at fixed composition against a zigzag and a clumped lithium
arrangement; and the $\sqrt3\times\sqrt3$ ordered structure at $x=1/3$, the composition toward
which the slowly delithiated batches move. Convergence settings, cell geometries, the magnetic
seeding protocol, and the caveats that attach to each number are in Sec.~S5 of
Ref.~\cite{SM}. Table~\ref{tab:dft} collects the results and Fig.~\ref{fig:dftu} displays them.

\begin{table}[!tbp]
  \caption{Density-functional results, at $U=\Uused$~eV on the cobalt $3d$ shell unless a row
  states otherwise. Magnetic energies are per cobalt relative to the ground state of the same
  cell, lithium-arrangement energies per formula unit relative to the row arrangement, and
  magnetic energy differences between alignments per Co$^{4+}$; no spin-Hamiltonian convention
  is used anywhere, so the last are plain energy differences and carry their signs. Every value
  in this table enters the text through a single macro defined in \texttt{dft\_numbers.tex}. The
  three blocks come from separate campaigns and are not directly comparable: the first is
  relaxed with the magnetic order converged, whereas the motif block is at fixed ideal geometry
  with ferromagnetic alignment imposed, which is why the row structure appears there with a
  slightly different gap. Section~S5 of Ref.~\cite{SM} carries the caveats, of which the most
  important is that $U$ is fixed by the requirement that the low-spin state be reproduced.}
  \label{tab:dft}
  \renewcommand{\arraystretch}{1.14}
  \begin{ruledtabular}
  \begin{tabular}{lr}
    Quantity & Value \\
    \colrule
    \multicolumn{2}{l}{\textit{Row-ordered} $x=1/2$, \textit{AFM ground state}}\\
    Co$^{4+}$ / Co$^{3+}$ local moment ($\muBs$) & \comomentRow{} / \comomentThreeRow \\
    $d$-occupation difference ($e$)              & \dnRow \\
    Co-O bond-length contrast (\AA)              & \bondRow \\
    Gap, AFM / FM (eV)                           & \gapRow{} / \gapRowFM \\
    $E_{\mathrm{FM}}-E_{\mathrm{AFM}}$ (meV/Co)  & $+$\efmafmRow \\
    $E_{\mathrm{NM}}-E_{\mathrm{AFM}}$ (meV/Co)  & $+$\enmRow \\
    $E_{\mathrm{AFM}}-E_{\mathrm{FM}}$, intra-row & \eafmIntra \\
    $E_{\mathrm{AFM}}-E_{\mathrm{FM}}$, inter-row & \eafmInter \\
    \quad both in meV per Co$^{4+}$              & \\
    \colrule
    \multicolumn{2}{l}{\textit{Lithium motifs at} $x=1/2$, \textit{fixed geometry, FM imposed}}\\
    $E_{\mathrm{zigzag}}-E_{\mathrm{row}}$       & \deltaEZigzag{} / \deltaEZigzagUzero \\
    $E_{\mathrm{clumped}}-E_{\mathrm{row}}$      & $+$\deltaEClumped{} / $+$\deltaEClumpedUzero \\
    \quad meV per f.u., $U=\Uused$~/~$0$         & \\
    Gap, row / zigzag / clumped (eV)             & \gapRowMotif/\gapZigzag/\gapClumped \\
    $D_{\mathrm{clumped}}/D_{\mathrm{row}}$ at $U=0$ & \dosContrast \\
    \colrule
    \multicolumn{2}{l}{$\sqrt3\times\sqrt3$ $x=1/3$, \textit{ferrimagnetic ground state}}\\
    Co$^{4+}$ local moment ($\muBs$)             & \comomentXthird \\
    $d$-occupation difference ($e$)              & \dnXthird \\
    Co-O bond-length contrast (\AA)              & \bondXthird \\
    Gap (eV)                                     & \gapXthird \\
    $E_{\mathrm{FM}}-E_{\mathrm{ferri}}$ (meV/Co) & $+$\efmafmXthird \\
    $E_{\mathrm{NM}}-E_{\mathrm{ferri}}$ (meV/Co) & $+$\enmXthird \\
    \quad the same per cell (meV)                & $+$\splitXthird \\
    $D_{\mathrm{NM}}(E_{F})$ (eV$^{-1}$/cell)    & \dosXthirdNM \\
  \end{tabular}
  \end{ruledtabular}
\end{table}

\begin{figure}[!tbp]
  \centering
  \figorpending{fig4_dftu}{\columnwidth}
  \caption{Density-functional results for the ordered phases, at $U=\Uused$~eV on the cobalt
  $3d$ shell. (a) Energies of the magnetic configurations, per cobalt and relative to the ground
  state of each composition. At $x=1/2$ the antiferromagnetic alignment within a Co$^{4+}$ row
  is lowest, the ferromagnetic alignment lies \efmafmRow~meV per cobalt above it, and the
  nonmagnetic solution lies \enmRow~meV above; at $x=1/3$ the ferrimagnetic arrangement is
  lowest, with the ferromagnetic one \efmafmXthird~meV and the nonmagnetic one \enmXthird~meV
  per cobalt above. (b) Energies of the three lithium motifs at $x=1/2$, per formula unit and
  relative to the row arrangement, filled markers at $U=\Uused$~eV and open markers at $U=0$.
  Clumping the stripes costs \deltaEClumped~meV per formula unit, above the reference line at
  $\kB\times240$~K${}=20.7$~meV, while the zigzag variant changes sign between the two values of
  $U$ and is therefore degenerate with the row arrangement. (c) Gap and Co$^{4+}$ moment for each
  phase, \gapRow~eV with \comomentRow~$\muBs$ at $x=1/2$ and \gapXthird~eV with
  \comomentXthird~$\muBs$ at $x=1/3$, against the nonmagnetic solutions, which have no gap at
  all; the panel prints the gaps to three decimals where the text rounds them to two. Moment, gap and charge order appear together, and removing the moment removes the other
  two.}
  \label{fig:dftu}
\end{figure}

The Hubbard parameter is calibrated rather than chosen. At $U=2.5$~eV the moment settles on the
lithium-rich cobalt instead of the lithium-poor one, inverting the charge order, and the
$x=1/3$ cell is metallic; at $U=5.5$~eV the Co$^{4+}$ site is driven high spin, with a moment
above $2\,\muBs$, and the $x=1/2$ cell is metallic again. Only near $U=\Uused$~eV does the
calculation reproduce the low-spin $S=1/2$ Co$^{4+}$ and $S=0$ Co$^{3+}$ configuration that the
chemistry of this material requires. That one requirement fixes the parameter, and every number
quoted below is quoted at that value and is meaningless without it. The Hubbard term remains a
stand-in for the screened exchange that a quasiparticle self-consistent $GW$ treatment with
electron-hole interactions supplies directly for the stoichiometric
parent~\cite{Radha2021c,Radha2023erratum}, so the gaps reported below label the charge-ordered
magnetic state and are not offered as predictions of a gap magnitude.

\subsection{Charge order, moment, and gap arrive together}

In every cell examined, the charge disproportionation, the local moment, and the gap appear
together or not at all. At $x=1/2$ the ground state is \magorderRow, carrying
\comomentRow~$\muBs$ on the lithium-poor cobalt against \comomentThreeRow~$\muBs$ on the
lithium-rich one, a $d$-occupation difference of \dnRow~electrons, and a gap of \gapRow~eV. The
moment is well determined, running from \comomentRowRange~$\muBs$ across the ferromagnetic and
both antiferromagnetic configurations, so nothing below turns on which of them is chosen. Aligning the same moments ferromagnetically costs \efmafmRow~meV per cobalt and
narrows the gap to \gapRowFM~eV, and removing them altogether costs \enmRow~meV per cobalt and
closes the gap: the nonmagnetic solution of the same cell is a metal. At $x=1/3$ the
$\sqrt3\times\sqrt3$ structure repeats the pattern, its ground state \magorderXthird, with the
cobalt beneath the lithium nonmagnetic and the other two carrying $\pm\,$\comomentXthird~$\muBs$, a $d$-occupation
difference of \dnXthird~electrons, and a gap of \gapXthird~eV, while its nonmagnetic solution is
a good metal with \dosXthirdNM~states per eV per cell at the Fermi level lying \enmXthird~meV
per cobalt higher. Forcing the moment to vanish closes the gap and erases the charge order at
one stroke. The phase the transport model calls O is therefore a magnetic semiconductor and not
a band insulator, and its three signatures are one transition seen from three directions.

The moment is the one the thermodynamics asked for. Values of \comomentRow{} and
\comomentXthird~$\muBs$ are the low-spin $S=1/2$ result for a $d^{5}$ cobalt in an octahedral
field, and the companion sites sit within \comomentThreeRow~$\muBs$ of zero, which is the $S=0$
$d^{6}$ value. Half the cobalt sites at $x=1/2$ thus carry a full spin one-half, so the ceiling of
\msatRow~$\muBs$ per cobalt that Eq.~(S35) uses to bound the latent entropy of a bulk
transformation is a computed quantity rather than an estimate.

The charge order is electronic and not structural. The site-mean Co-O bond lengths of the two
cobalt sites differ by \bondRow~\AA{} at $x=1/2$ and \bondXthird~\AA{} at $x=1/3$ while their
$d$~occupations differ by \dnRow{} and \dnXthird~electrons. Individual bonds are not that
uniform, spanning \bondSpanRow~\AA{} within a single octahedron, so what is nearly absent is the
contrast between the two cobalt sites rather than displacement in general; a displacive or
Jahn-Teller transition carrying the same electronic contrast would separate the two sites by
tens of picometers. That difference is measurable, and it becomes the spectroscopic
entry of Table~S2 in Ref.~\cite{SM}: diffraction and phonon spectroscopy should see almost nothing
across this transformation while tunneling spectroscopy sees a hard gap open.

\subsection{The energy scale of the ordering}

Breaking the lithium stripes into clumps costs \deltaEClumped~meV per formula unit at
$U=\Uused$~eV and \deltaEClumpedUzero~meV per formula unit at $U=0$, both with ferromagnetic
alignment imposed and at fixed ideal geometry. That is the same order as the thermal scale of
the collapse, $\kB$ times 240~K being 20.7~meV, so a configurational excitation of a few tens of
meV per formula unit places the order-disorder boundary in the few-hundred-kelvin range with
nothing adjusted, and it is this energy that supplies the driving force $\Delta f$ of
Eq.~\eqref{eq:kjmarates} and the nucleation barrier of Eq.~(S21). The clumped pattern is a
clustering proxy and not a disordered state, so the comparison fixes an order of magnitude
rather than a transition temperature.

Rerouting a stripe, by contrast, costs almost nothing at all. The zigzag arrangement lies \deltaEZigzag~meV per
formula unit from the row arrangement at $U=\Uused$~eV and \deltaEZigzagUzero~meV per formula
unit at $U=0$, so the sign reverses with the Hubbard parameter and the two motifs are degenerate
within the resolution of the method. The ordered phase is accordingly not one arrangement but a
manifold of nearly degenerate stripe variants, which is what makes it domain textured. That
degeneracy is not an entropy reservoir: $\kB\ln3$ per domain amounts to $2\times10^{-4}\,\kB$
per cobalt for domains 20~nm across, four hundred times below the entropy scale the field
measurement works with. What the field measurement
states is the separate and stronger proposition that the transformation converts only part of
the material, its latent entropy bounded by Eq.~(S35) at $\dSbound\,\kB$ per cobalt under the
bulk reading of Sec.~\ref{sec:field}, which is about a third of the $\kB\ln2$ that complete
lithium order and disorder at half filling would release.

\subsection{What the calculations exclude}

There are two things the mechanism might have been and is not.

The first is that the resistive contrast comes from which periodic lithium motif is present, and
the calculations do not support it. The comparison has to be made in the quantity that controls
conduction, and since transport here is not activated across a gap, which Sec.~\ref{sec:mixture}
establishes from the size of the measured excess alone, that quantity is the density of states
at the Fermi level. At $U=0$, where all three arrangements are metals, their densities of states
at the Fermi level agree to within one percent, the clumped-to-row ratio being \dosContrast{} on
an absolute value of \dosRowUzero~states per eV per cell, so the three motifs are
indistinguishable in the one number that could have produced $s>1$. At $U=\Uused$~eV all three
are instead charge-ordered insulators, with gaps of \gapRowMotif, \gapZigzag{} and
\gapClumped~eV for the row, zigzag and clumped patterns and a density of states at the Fermi
level that is numerically zero in each; that spread of gaps is not a transport statement, since
a gap-activated channel would give the contrasts of $10^{7}$ and larger that the measurement
excludes, and what the three insulating solutions establish instead is that the charge-ordered
character of the ordered phase is itself motif independent. Independent evidence points the same
way, bulk Li$_{0.5}$CoO$_2$ remaining metallic through its own ordering transition, its in-plane
resistivity staying between 0.75 and 0.95~m$\Omega$~cm from 4.2 to 300~K with no metal-insulator
transition anywhere~\cite{OuYang2012}. Whatever makes the ordered phase more resistive, it is not
the lithium pattern acting on the cobalt bands, because the cobalt charge order forms equally
well on top of any of these patterns and is indifferent to which. The factor $s>1$ in
Eq.~\eqref{eq:mixture} cannot be assigned to a motif.

A different mechanism survives, and these calculations do not test it. All three arrangements
are periodic, so a commensurate cobalt charge order is available in every one of them, whereas
genuine lithium disorder is not representable in a cell this size. The natural reading is that
in the kinetically frozen phase the aperiodic lithium occupation frustrates the commensurate
cobalt charge order, leaving weight inside the gap and a conductivity higher than the cleanly
ordered phase achieves, so that the transformation raises the resistance by removing frustration
rather than by opening a gap where none existed. That reading is a hypothesis and not a
result, and the price of routing the exclusion through the metallic densities of states is that
the calculations then supply no estimate of $s$ at all: the contrast of 2 to 5 is read off the
measurement, and the fifth to a half of the low-energy weight that it says the ordered phase
retains is the target the quasirandom calculation has to hit. The calculation that decides it is
a special quasirandom structure large enough to carry an aperiodic lithium occupation at fixed
$x=1/2$, with its density of states at the Fermi level compared against the row-ordered cell at
the same $U$; the $4\times2$ cell used here cannot. What the gap the calculations do return is
good for is labelling the charge-ordered magnetic state, and not for setting a transport scale.

The second exclusion concerns the field, and it turns on a sign rather than on a magnitude. We
report the magnetic energetics as plain energy differences per Co$^{4+}$, so that no
spin-Hamiltonian convention is smuggled in. Within a Co$^{4+}$ row at $x=1/2$ the
antiferromagnetic alignment lies below the ferromagnetic one,
$E_{\mathrm{AFM}}-E_{\mathrm{FM}}=\eafmIntra$~meV per Co$^{4+}$; between rows the sign reverses,
\eafmInter~meV per Co$^{4+}$; and at $x=1/3$ the ferrimagnetic ground state lies \splitXthird~meV
per cell below the ferromagnetic one. The dominant coupling is antiferromagnetic at both
compositions, and both computed ground states therefore carry zero net moment.

A large exchange is not in itself an obstacle to a field response; it is the very thing that
would hold a cluster of twenty moments rigidly parallel so that a weak field could turn the
cluster whole, and the ferromagnetically correlated Co$^{4+}$ domain of Sec.~\ref{sec:field}
would be helped by it. What excludes the bulk sublattice is that its moments are not parallel:
Eq.~\eqref{eq:cc} asks the ordered phase for more magnetization than the parent, and an
antiferromagnet offers none. The magnitude adds only the separate statement that 5~T, whose
Zeeman energy is 0.58~meV against couplings near \jintraMag~meV, cannot rearrange that spin
structure into one that would.

A ferromagnetically polarized reservoir is therefore required, and the most natural candidate in
this material is the spin-polarized two-dimensional electron gas of CoO$_2$-terminated
surfaces~\cite{Radha2021}, polarized by construction and resident on the surfaces of a flake
thin enough for surfaces to carry a large share of the conductance. It is the best of the
candidates considered here rather than the only one the material admits: oxygen $2p$ holes,
which delithiation is known to create~\cite{Salagre2023}, and moments on lithium-vacancy
clusters, domain walls or flake edges have not been examined here, and any of them could in
principle supply a polarized cluster. What distinguishes a surface reservoir from all of them is
its thickness dependence, which Sec.~\ref{sec:tests} turns into the measurement that decides the
question.

\section{Predictions and discriminating measurements}
\label{sec:tests}

The model as it stands reproduces four numbers read from published traces, the onset, the
collapse, the amplitude, and the field ratio, using four parameters, the barrier, the collapse
temperature, the product $c\phi(s-1)$, and the magnetization difference. That is a
parameterization and not yet a theory; it respects every
constraint, which no previous description of these data did, and it becomes a theory only
through predictions that were not used in building it.
Table~S2 of Ref.~\cite{SM} collects them, each set beside the measurement that decides it. Every
entry is out of sample, and every entry can be performed on existing devices with existing
instrumentation.

Three of them carry most of the weight. The warming-rate series tests the central claim
that the onset is a clock reading rather than a phase boundary, and Eq.~\eqref{eq:rateshift}
predicts its slope with no freedom beyond the barrier already fixed by the onset itself; a flat
onset across a factor of fifty in rate, at a cooling protocol held fixed throughout, would end
the kinetic account outright. The warming-branch hold tests the growth-from-seeds picture in the
most direct way available, by giving the transformation the one thing a sweep denies it, which
is time at a fixed temperature. The thickness series has been promoted
by the first-principles results of Sec.~\ref{sec:dft}, which leave the surface electron gas as
the only remaining reservoir of field-coupled moments: if the excess sheet conductance is
independent of thickness, the ordered phase and its moments live on the surfaces, and if it
scales with thickness, the field response needs an explanation this paper does not have.

The other entries each remove one surviving alternative: a cooling-rate series at fixed warming
rate settles the provenance of the seeds, which is the one ingredient of
Sec.~\ref{sec:asymmetry} assumed rather than derived, since athermal templates leave the onset
where it is while thermally born ones raise it by 5.8~K per decade; a field series at 1, 2, 3,
5, 7 and 9~T, read together with a single in-plane run, separates the three couplings of
Sec.~\ref{sec:field} by the shape of the amplitude alone; channel-length scaling, a four-probe
measurement and a bias-history series close the contact-artifact alternative of
Sec.~\ref{sec:constraints}, including the Debye-thin space-charge layer that the
blocking-contact argument leaves standing; calorimetry across the anomaly, read together with
the thickness series, fixes which cobalt population Eq.~\eqref{eq:dM} refers to, a bulk
transformation carrying 2.1~J/g against the 0.02~J/g of a surface one; a batch series in $x$ at
fixed protocol tests whether the onset moves only through the composition dependence of $\Ea$,
so that a shift beyond about 20~K, or one uncorrelated with the width of the rise, implicates
inhomogeneity or a second transition rather than a single kinetic edge; and electron or
synchrotron diffraction on one flake at 100~K and again at 300~K, with weak lithium-contrast
superlattice intensity expected in the warmed-through low-temperature state and absent in the
as-cooled state at an interlayer spacing that does not move, confirms that what orders is the
lithium sublattice and excludes the stacking glide of Sec.~\ref{sec:lisublattice}.

\section{Discussion}
\label{sec:discussion}

Three features of this account deserve emphasis because they are what distinguish it from the
description it replaces.

The first is that the onset temperature carries no information about any electronic energy
scale. It is a solution of Eq.~\eqref{eq:Ea}, a transcendental relation among a barrier, an
attempt frequency, a sweep rate and a seed density, and it moves when any of the four moves, so
matching such a number against an equilibrium crossover computed for an electronic model is a
category error independent of how close the two numbers happen to fall. The bulk material offers
the same test on better material. Its 155 to 185~K anomaly is first order, hysteretic and
latent-heat bearing, and the rate dependence reported for those crystals belongs to the separate
and lower lithium-freezing scale near 130~K~\cite{Miyoshi2010,Motohashi2009}, so the bulk
anomaly is not itself an unfreezing. What the present model predicts for it is specific: a
warming-rate series on bulk crystals should move a kinetic edge near 155~K by about 12~K per
decade while a thermodynamic edge stays fixed to better than $10^{-4}$~K, which decides the
character of each bulk scale on its own terms. The mapping carries a cost, in that the causal
order proposed for the bulk runs opposite to the one used here, cobalt charge order
below the upper scale templating the lithium order below the lower~\cite{Miyoshi2010}, whereas
this account has the lithium sublattice ordering first and the cobalt charge order following.
What we offer against it is the finding of Sec.~\ref{sec:dft} that the cobalt charge order forms
equally well on every lithium motif examined and so cannot be slaved to the lithium pattern,
which reads more naturally as a cooperative transition than as one order templating the other;
the same bulk warming-rate series decides that too.

The second is that the field argument has become sharper by losing a candidate, and that it
remains conditional on the sign of a shift the present data do not establish. It asks only that
the ordered phase carry a magnetization the parent phase does not, which Eq.~\eqref{eq:cc}
converts into a tilt of the collapse, and that the moments be correlated over blocks of order
ten, which Eq.~\eqref{eq:headcount} obtains from polarization arithmetic alone. The
first-principles calculations supply the moment, \comomentRow~$\muBs$ of low-spin $S=1/2$ on
half the cobalt sites, and then take away the obvious place to put it, since within the bulk
those moments prefer to be antiparallel and an antiferromagnet presents nothing for a field to
grip. A strong exchange would help rather than hinder, holding a block of moments rigid enough
for a weak field to turn it whole, so it is the sign of the coupling and not its size that does
the excluding. A ferromagnetically polarized reservoir is therefore required if the shift is
real, and the spin-polarized surface electron gas~\cite{Radha2021} is the candidate named in
Sec.~\ref{sec:dft}, a polarized sheet of correlated moments on exactly the surfaces that
dominate a flake tens of nanometers thick.

The third is a failure that changes what the model claims, and it is reported here in full. The
calculations were meant to show that the lithium-ordered arrangement has the smaller density of
states at the Fermi level, and they show no such thing: at $U=0$, where the comparison bears on
transport at all, the row, zigzag and clumped arrangements are metals whose densities of states
at the Fermi level agree to within one percent, and at $U=\Uused$~eV all three are instead
charge-ordered insulators with gaps within \gapMotifSpread~eV of one another. The cobalt charge
order forms on top of any lithium pattern and does not care which. What the transport model
needs, therefore, is not lithium order but the difference between order and genuine disorder,
and the mechanism left standing is that aperiodic lithium occupancy frustrates the commensurate
cobalt charge order and leaves conducting weight in the gap, in the amount, a fifth to a half of
the weight the frozen phase carries, that the measured contrast requires. That is a testable
claim about a calculation, not a rhetorical retreat: a special quasirandom supercell at $x=1/2$
either shows in-gap weight relative to the row-ordered cell or it does not.

Scanning tunneling spectroscopy of this surface finds the ordered region conducting and the
structurally disordered region gapped~\cite{Iwaya2013}, whereas the transformation analyzed here
makes the ordering raise the resistance. These are different orders and different observables,
a surface cobalt registry probed by local density of states in the first case and volume lithium
and charge order probed by two-terminal transport in the second. They are not in contradiction,
and the frustration mechanism proposed above in fact predicts a disordered region with weight
inside the gap, but only the thickness series decides which of the two the flake anomaly
reports.

Beyond this material, the argument is a reminder that cathode oxides are routinely cooled
through their ionic ordering temperatures at rates nobody records, and that the calorimetric
tools built for glasses apply to them without modification. Nothing about the readout is peculiar
to this system either, a structural transition of a mobile cation sublattice having already been
followed through the two-terminal conductance of a single Ag$_2$Te
nanowire~\cite{Premasiri2019}. A hysteresis loop whose warming
branch differs from its cooling branch is not a curiosity to be noted and set aside; it is a
measurement of a barrier, provided one records the sweep rate.

\section{Conclusions}
\label{sec:conclusions}

The warming-only resistance anomaly of slowly delithiated Li$_x$CoO$_2$ nanoflakes follows from
cold crystallization of the lithium sublattice. Classical nucleation and growth are never
favorable together along the cooling sweep, so the flake reaches base temperature as an
untransformed matrix carrying a quenched population of defect-templated embryos, which the
model adopts as an assumption and which a cooling-rate series can confirm, and the warming sweep
grows them into a more resistive ordered phase as soon as lithium moves. The unfreezing
condition fixes an effective barrier of $0.40\pm0.05$~eV, together with a seed density between
$10^{-4}$ and $10^{-1}$~nm$^{-2}$, which identifies the divacancy hop as the mechanism and
predicts an onset shift of $12\pm1$~K per decade of warming rate at fixed cooling protocol,
independent of the Avrami exponent. A field-shifted collapse would require the ordered phase to
carry the larger magnetization, bounding the latent entropy below about $\dSbound\,\kB$ per Co
if the transforming phase is the bulk and by two orders of magnitude less if it is the surfaces,
and the polarization arithmetic would require the field-coupled unit to hold of order ten
correlated moments; the published shift is consistent with zero at one standard deviation and
our digitization of the same traces returns the same sign at two standard deviations in the raw
branch difference and at one in the normalized excess, so that whole chain waits on a field
series.

Density-functional calculations with a Hubbard $U$ fixed by the low-spin requirement confirm the
ingredients and remove two readings. Charge order, a Co$^{4+}$ moment of \comomentRow~$\muBs$,
and a gap of \gapRow~eV at $x=1/2$ and \gapXthird~eV at $x=1/3$ appear and vanish together, the
nonmagnetic solutions being metals; clumping the lithium stripes costs \deltaEClumped~meV per
formula unit, the same order as the thermal scale of the collapse, while rerouting them costs
nothing measurable, so the transformation is a partial ordering within a nearly degenerate
stripe manifold; and the charge order contrasts the site-mean Co-O bond lengths by only
\bondRow{} to \bondXthird~\AA, so it is electronic and nearly invisible to diffraction. Against
this, the three periodic lithium motifs are indistinguishable in the density of states at the
Fermi level that controls conduction, so the resistive contrast is not a matter of which motif
is present, and the surviving candidate is frustration of the cobalt charge order by aperiodic
lithium, which a large quasirandom supercell can test and this work has not. The bulk Co$^{4+}$
moments are coupled antiferromagnetically and so present no net magnetization for a field to act
on, which leaves a ferromagnetically polarized reservoir as a requirement and the surface
electron gas of CoO$_2$-terminated Li$_x$CoO$_2$ as its most natural realization among the
candidates considered. A warming-rate series at fixed cooling protocol, a cooling-rate series, a
hold on the warming branch, a field series, and a thickness series read together with
calorimetry decide the picture with measurements that need no new instrumentation.

\section*{Data availability}

The scripts that generate every figure, the traces digitized from the published figures of
Refs.~\cite{Crowley2020,RadhaPerspective2021}, and the \textsc{Quantum ESPRESSO} input files,
structure files and distilled outputs behind every number in Table~\ref{tab:dft} are available
from the author and are deposited in the repository accompanying this manuscript. The raw
resistance traces underlying Refs.~\cite{Crowley2020,RadhaPerspective2021} are held by the
originating collaboration.

\begin{acknowledgments}
The author thanks [acknowledgments to be completed] for discussions. This work was supported by
[funding to be completed]. Computational resources were provided by [computing allocation to be
completed].
\end{acknowledgments}

\bibliographystyle{apsrev4-2}
\bibliography{refs}

\end{document}
